
This paper studies the decidability and complexity of the existential
theory of an extension of the structure
$\langle \N,0,1,+,2^x \rangle$, where $2^x$ is the function mapping a
natural number $n$ to $2^n$.  Decidability of the first-order theory of
this structure was first shown by Sem{\"e}nov in a more general
framework using an automata-theoretic approach~\cite{Sem80}, and we
henceforth call this theory
\emph{Sem{\"e}nov arithmetic}. As shown by Cherlin and Point~\cite{CP86},
see also~\cite{Point10}, Sem{\"e}nov arithmetic admits quantifier
elimination and has a quantifier-elimination procedure that runs in
non-elementary time, and this upper bound is tight~\cite{Point10}. The
existential fragment of Sem{\"e}nov arithmetic has recently been shown
decidable in NEXP~\cite{BCM23} by giving a more elaborate quantifier
elimination procedure. Unlike its substructure Presburger arithmetic
(obtained from dropping the function $2^x$), Sem{\"e}nov arithmetic is
not automatic in the sense of the theory of automatic
structures~\cite{KN95,BG00}. The constant growth lemma~\cite{KM07}
states that for any function $f$ of an automatic structure,
$\abs{f(x_1,\ldots,x_n)} \le \abs{x_1} + \cdots + \abs{x_n} + c$ for
some constant $c$. This is clearly not the case for the function
$2^x$.

The decidability of Sem{\"e}nov arithmetic is fragile with respect to
extensions of the structure. For instance, it is not difficult to see
that extending Sem{\"e}nov arithmetic with the B\"uchi predicate
$V_2(x,y)$, where $V_2(x,y)$ holds whenever $x$ is the largest power
of two dividing $y$ without remainder, results in an undecidable
first-order theory, see e.g.~\cite{Point10}. However, this
undecidability result requires an $\exists^*\forall^*$-quantifier
prefix and does not rule out decidability of the existential
fragment. The main result of this paper is to show that the
existential theory of \emph{generalised Sem{\"e}nov arithmetic}, i.e.,
the existential theory of $\langle
\N,0,1,+,2^x,\{R_i\}_{i\ge 0} \rangle$, where $R_0,R_1,\ldots$ is an
enumeration of all regular languages over the alphabets $\{0,1\}^d$,
$d\ge 1$, is decidable in EXPSPACE. Non-automaticity of Sem{\"e}nov
arithmetic and undecidability of $\langle \N,0,1,+,2^x,V_2 \rangle$
rule out the possibility of approaching this existential theory via
automatic structures based on finite-state automata or via
quantifier-elimination \emph{\`a la} Cherlin and Point, since $V_2$ is
definable as a regular language over pairs of number strings.
Instead, our decidability result is based on a reduction to the
language emptiness problem of a special class of \emph{affine vector
addition systems with states (affine VASS)}.

A VASS comprises a finite-state controller with a finite number of
counters ranging over the natural numbers. In an affine VASS, when
taking a transition, every counter can be updated by applying an
affine function $x \mapsto a x + b$ to the current value,
provided that the resulting counter is non-negative. While
reachability in affine VASS is decidable for a single
counter~\cite{FGH13}, already in the presence of two counters
reachability becomes undecidable~\cite{Reichert2015}. Our reduction
consequently requires a restricted class of affine VASS to obtain
decidability. We call this class
\emph{restricted labelled affine VASS (restricted \laVASS)}. A
restricted \laVASS is an affine VASS with $d$ pairs of counters and
hence $2d$ counters in total. For every pair, the first counter does
not change until it keeps getting incremented at every transition; the
second counter is only updated via affine functions $x\mapsto 2
x$ and $x\mapsto 2 x+1$. A configuration consisting of a control
state and $2d$ counter values is accepting whenever the control state
is accepting and for every pair of counters, the first counter has the
same value as the second counter. We give an EXPSPACE procedure for
deciding emptiness of restricted \laVASS whose correctness proof is
based on a kind of counter elimination procedure in which we
successively encode counters into a finite state space while
preserving equi-non-emptiness. The tight syntactical restrictions
on \laVASS are necessary in order to obtain a decidable class of
affine VASS---relaxing those restrictions even only slightly leads to
undecidability of the language emptiness problem as we will later
discuss.

The EXPSPACE upper bound for existential generalised Sem{\"e}nov
arithmetic follows from a reduction to language non-emptiness of a
restricted \laVASS whose language encodes all solutions of the given
formula. Obtaining an elementary upper bound is difficult since it is
easily seen that smallest solutions of an existential formula of
Sem{\"e}nov arithmetic can be non-elementary in bit-length.

As an application of our EXPSPACE upper bound for existential
generalised Sem{\"e}nov arithmetic, we show that a certain class of
string constraints with length constraints is decidable in
EXPSPACE. It allows existentially quantifying over bit-strings, and to
assert that the value of a string variable lies in a regular language,
as well as Presburger-definable constraints over the lengths of the
bit-strings stored in string variables and the numerical values of
those variables (when viewed as encoding a number in binary).
Decidability of this class was left open in~\cite{BDGKMMN23}. We
settle this open problem by showing that it can be reduced to the
existential fragment of generalised Sem{\"e}nov arithmetic. Formulas
of this class of string constraints appear widely in practice---in
fact, essentially all formulas in the extensive collection of standard
real-world benchmark sets featured in~\cite{BDGKMMN23,BKMMDNG20} lie
in this class.
